% Part: turing-machines
% Chapter: undecidability
% Section: enumerating-tms

\documentclass[../../../include/open-logic-section]{subfiles}

\begin{document}

\olfileid{tur}{und}{enu}
\olsection{Ngenumerasi Mesin Turing}

\begin{explain}
Kita bisa nuduhake yen himpunan wakil mesin Turing standar nganti panggantos jeneng iku
!!{enumerable}. Cathetan koreksi sumber OLPL-772: iki mung bener kanggo wakil mesin standar nganti panggantos jeneng kahanan lan simbol; mesin mentah kanthi label obyek sembarang bisa ora enumerable. Iki ndherek saka kasunyatan yen saben
mesin Turing bisa diterangake kanthi winates. Himpunan
kahanan lan kosakata pita iku himpunan winates. Fungsi
transisi iku fungsi parsial saka $Q \times \Sigma$
menyang $Q \times \Sigma \times \{\TMleft, \TMright, \TMstay\}$,
mula uga bisa diterangake kanthi ndhaptar nilaine
kanggo pasangan argumen winates cacahe kang dadi
daerah tegesé.

Nganti kene pratelan iku bener, nanging ana prabédan
kang alus. Definisi mesin Turing ora matesi apa wae
kang bisa dadi !!{element}s ing himpunan kahanan
lan alfabet pita. Mula, contone, kanggo saben
wilangan real, sacara teknis ana mesin Turing
kang nggunakake wilangan iku minangka kahanan.
Nanging \emph{tumindake} mesin Turing ora gumantung
marang obyèk apa kang dadi kahanan lan simbol.
Gatekna rong mesin Turing ing \olref{fig:variants}.
\begin{figure}
\begin{center}
\begin{tikzpicture}[->,>=stealth',shorten >=1pt,auto,node distance=2.8cm,
                    semithick]
  \tikzstyle{every state}=[fill=none,draw=black,text=black]

  \node[initial,state]         (A)              {$q_0$};
  \node[state]         (B) [right of=A] {$q_1$};

  \path (A) edge [bend left] node {\TMtrans{\TMstroke}{\TMstroke}{\TMright}} (B)
        (B) edge [loop above] node {\TMtrans{\TMblank}{\TMblank}{\TMright}} (B)
            edge [bend left] node {\TMtrans{\TMstroke}{\TMstroke}{\TMright}} (A);
\end{tikzpicture}\\
\begin{tikzpicture}[->,>=stealth',shorten >=1pt,auto,node distance=2.8cm,
  semithick]
\tikzstyle{every state}=[fill=none,draw=black,text=black]

\node[initial,state]         (A)              {$s$};
\node[state]         (B) [right of=A] {$h$};

\path (A) edge [bend left] node {\TMtrans{A}{A}{\TMright}} (B)
(B) edge [loop above] node {\TMtrans{\TMblank}{\TMblank}{\TMright}} (B)
edge [bend left] node {\TMtrans{A}{A}{\TMright}} (A);
\end{tikzpicture}
\end{center}
\caption{Varian mesin \emph{Genep}}
\ollabel{fig:variants}
\end{figure}
Rong diagram iki nggambarake rong mesin: $M$ kanthi
alfabet pita $\Sigma = \{\TMendtape,\TMblank,\TMstroke\}$
lan himpunan kahanan $\{q_0,q_1\}$, sarta $M'$ kanthi
alfabet $\Sigma' = \{\TMendtape,\TMblank,A\}$ lan
kahanan $\{s,h\}$. Nanging instruksi liyane padha:
$M$ mandheg nalika maca rerangken $n$~$\TMstroke$
yen lan mung yen $n$ genep, lan $M'$ mandheg nalika
maca rerangken $n$~$A$ yen lan mung yen $n$ genep.
Kita mung ngganti jeneng $\TMstroke$ dadi~$A$,
$q_0$ dadi~$s$, lan $q_1$ dadi~$h$. Tuladha iki
bisa digeneralisasi: kita bisa nganggep rong mesin
Turing padha yen siji bisa diasilake saka sijine
kanthi ngganti jeneng simbol lan kahanan kaya iki.
Malah simbol lan kahanan mesin Turing bisa wae
kita anggep minangka wilangan wutuh positif:
tinimbang $\sigma_0$ gunakna~$1$, tinimbang
$\sigma_1$ gunakna~$2$, lan sateruse;
$\TMendtape$ dadi~$1$, $\TMblank$ dadi~$2$,
lan sateruse. Kanthi cara iki, mesin \emph{Genep}
dadi mesin kang digambarake ing
\olref{fig:standard-even}.
\begin{figure}
\[\begin{tikzpicture}[->,>=stealth',shorten >=1pt,auto,node distance=2.8cm,
  semithick]
\tikzstyle{every state}=[fill=none,draw=black,text=black]

\node[initial,state]         (A)              {$1$};
\node[state]         (B) [right of=A] {$2$};

\path (A) edge [bend left] node {\TMtrans{3}{3}{\TMright}} (B)
(B) edge [loop above] node {\TMtrans{2}{2}{\TMright}} (B)
edge [bend left] node {\TMtrans{3}{3}{\TMright}} (A);
\end{tikzpicture}
\]
\caption{Mesin \emph{Genep} standar}
\ollabel{fig:standard-even}
\end{figure}
Mesin Turing kang kahanan lan simbole awujud
wilangan wutuh positif bisa kita arani mesin
\emph{standar}, lan wiwit saiki kita mung bakal
nggatekake mesin standar.\footnote{Istilah
``mesin standar'' iku dudu istilah baku.}

Kita kepengin nuduhake yen himpunan wakil mesin Turing iku
!!{enumerable}; kanthi tetimbangan ing ndhuwur, cukup
nuduhake yen himpunan mesin Turing standar iku
!!{enumerable}. Upamane ana mesin Turing standar
$M = \tuple{Q, \Sigma, q_0, \delta}$. Kepriye carane
njlentrehake mesin iku nganggo rerangken winates
wilangan wutuh positif? Dhisik kita dhaptar cacahe
kahanan, kahanan-kahanane, cacahe simbol, simbol-simbole,
lan kahanan wiwitan. (Elinga, kabeh iki wilangan
wutuh positif amarga $M$ mesin standar.) Banjur
kepriye karo~$\delta$? Himpunan argumen kang mungkin,
yaiku pasangan $\tuple{q,\sigma}$, winates amarga
$Q$ lan~$\Sigma$ winates. Mula katrangan ngenani
$\delta$ mung dhaptar winates kabeh tupel-$5$
$\tuple{q, \sigma, q', \sigma', d}$ kang
$\delta(q,\sigma) = \tuple{q', \sigma', D}$, lan
$d$ wilangan kang ngode arah~$D$ (upamane,
$1$ kanggo~$\TMleft$, $2$ kanggo~$\TMright$,
lan $3$ kanggo~$\TMstay$).

Kanthi cara iki, saben mesin Turing standar bisa
diterangake nganggo dhaptar winates wilangan
wutuh positif, yaiku minangka rerangken $s_M \in
(\PosInt)^*$. Contone, mesin \emph{Genep} standar
dikode nganggo rerangken iki:
\[
2, \underbrace{1, 2}_Q, 3, \overbrace{1, 2, 3}^\Sigma, 1, \underbrace{1, 3, 2, 3, 2}_{\delta(1,3) = \tuple{2,3,R}}, 
\overbrace{2, 2, 2, 2, 2}^{\delta(2,2) = \tuple{2,2,R}},
\underbrace{2, 3, 1, 3, 2}_{\delta(2,3) = \tuple{1,3,R}}.
\]
\end{explain}

\begin{thm}
Ana fungsi saka $\Nat$ menyang~$\Nat$ kang
ora bisa diitung dening mesin Turing.
\end{thm}

\begin{proof}
Kita ngerti yen himpunan rerangken winates wilangan
wutuh positif~$(\PosInt)^*$ iku !!{enumerable}
(\cref{sfr:siz:zigzag:prob:posint-star}). Mula
himpunan andharan mesin Turing standar, minangka
subhimpunan saka~$(\PosInt)^*$, uga bisa
dienumerasi. Saben fungsi saka $\Nat$ menyang~$\Nat$
kang bisa diitung dening mesin Turing diitung
dening sawijining mesin Turing (malah akeh mesin).
Kanthi ngganti jeneng kahanan lan simbol dadi
wilangan wutuh positif (mligine $\TMendtape$
dadi~$1$, $\TMblank$ dadi~$2$, lan $\TMstroke$
dadi~$3$), kita weruh yen saben fungsi kang
bisa diitung dening mesin Turing uga bisa
diitung dening mesin Turing standar. Tegese,
himpunan kabeh fungsi saka $\Nat$ menyang~$\Nat$
kang bisa diitung dening mesin Turing uga bisa
dienumerasi. Cathetan panyunting: enumerable ing bukti iki teges kardinalitas, ora pratelan ana prosedur efektif kanggo nyaring kabeh mesin kang ngitung fungsi total.

Nanging, himpunan kabeh fungsi saka $\Nat$
menyang~$\Nat$ ora !!{enumerable}
(\cref{sfr:siz:red:prob:nat-nat}). Yen kabeh
fungsi bisa diitung dening mesin Turing,
kita mesthi bisa ngenumerasi himpunan kabeh
fungsi kanthi ndhaptar kabeh andharan mesin
Turing kang ngitung fungsi-fungsi iku. Mula
ana fungsi kang ora bisa diitung dening
mesin Turing.
\end{proof}

\begin{prob}
  Apa kowe bisa nemokake cara njlentrehake mesin
  Turing tanpa kudu ndhaptar kahanan lan simbol
  alfabet kanthi eksplisit? Kowe oleh netepake
  pangertenmu dhewe bab mesin ``standar'', nanging
  % OLPL-150: Sumber beku nyebut mesin Turing diitung dening mesin standar; maksud soal yaiku tumindake bisa ditiru sawise panggantos representasi.
  Cathetan koreksi sumber OLPL-150: maksude yaiku niru tumindak mesin, dudu ngitung mesin minangka asil fungsi. Terangna apa sebabe tumindake saben mesin Turing
  bisa ditiru dening mesin ``standar'' miturut
  pangerten anyar iku.
\end{prob}

\end{document}
